\documentclass[utf8]{FrontiersinHarvard}
\usepackage{url,hyperref,lineno,microtype,subcaption,amsmath,booktabs,multirow}
\usepackage[onehalfspacing]{setspace}
\linenumbers

\def\keyFont{\fontsize{8}{11}\helveticabold }
\def\firstAuthorLast{Shen, Lei, Hung, Li, Pan and Fu {et~al.}}
\def\Authors{Tongfei Shen\,$^{1,2\dagger}$, Xueqin Lei\,$^{3\dagger}$, Hung Ka Lok\,$^{4\dagger}$, Huaicheng Li\,$^{5}$, Zhongze Pan\,$^{6}$ and Xuepeng Fu\,$^{7,*}$}
\def\Address{$^{1}$School of Information Science and Engineering, Qingdao Huanghai University, Qingdao, China \\
$^{2}$Department of Computer Science, City University of Hong Kong, Kowloon, China\\
$^{3}$College of Physics and Electronic Information Engineering, Zhejiang Normal University, Jinhua 321004, China\\
$^{4}$The Centre for Innovation and Entrepreneurship, The Hang Seng University of Hong Kong, Hong Kong, China\\
$^{5}$Ability R\&D Energy Research Centre, School of Energy and Environment, City University of Hong Kong, Kowloon, Hong Kong, China\\
$^{6}$School of Engineering, Nanfang College Guangzhou, Guangzhou, China\\
$^{7}$Department of Life Science and Agroforestry, Qiqihar University, Qiqihar 161000, China}
\def\corrAuthor{Xuepeng Fu}
\def\corrAddress{Department of Life Science and Agroforestry, Qiqihar University, Qiqihar 161000, China}
\def\corrEmail{02383@qqhru.edu.cn}

\begin{document}
\onecolumn
\firstpage{1}

\title[SIMPLEX: composition-space deep learning for hydrogel adhesion screening]{SIMPLEX: Composition-Space Deep Learning for Bio-inspired Hydrogel Adhesion with Prospective Validation on Model-Discovered Formulations}

\author[\firstAuthorLast ]{\Authors}
\address{}
\correspondance{}
\extraAuth{}
\maketitle

\begin{abstract}
\section{}
The design of functional hydrogels for biohybrid interfaces -- wet, load-bearing junctions between living tissue and synthetic materials -- relies on trial-and-error experimentation, because every candidate formulation must be synthesised and mechanically characterised before its underwater adhesion can be assessed. A data-driven model that predicts adhesion strength from monomer composition could accelerate candidate screening in this data-scarce regime. Here we introduce SIMPLEX, a dual-modality residual network with interaction self-attention that maps six monomer molar fractions (plus their 15 pairwise interaction terms, encoding composition synergy explicitly) to underwater glass adhesion strength (kPa), trained under Mixup, stochastic weight averaging and a range-domain constraint. We train on 316 hydrogel formulations spanning the full experimentally explored adhesion range (1 to 353 kPa; 19\% above 150 kPa) and prospectively validate on 25 held-out formulations discovered in the final round of the original study's sequential model-based optimisation loop (adhesion 62 to 251 kPa) -- a held-out, model-guided cohort that was never touched during model selection. Across 5-fold grouped cross-validation repeated over 5 seeds, SIMPLEX reaches $R^2 = 0.79$, statistically tied with the strongest baseline (random forest, $0.81$). On the 25-formulation prospective cohort, SIMPLEX attains $R^2 = 0.71$, Spearman $\rho = 0.87$, a tier-level ROC-AUC of 0.94, and Top-10/Top-20 screening precision of 1.00/0.90 -- a level shared by all baselines, indicating that the prospective ranking task is highly predictable; SIMPLEX differs in attaining the numerically highest $R^2$ ($0.71$) and Spearman correlation ($0.87$) among equally tuned baselines (random forest $0.56$, SVR $0.63$, Ridge $0.64$), although the cohort is too small to separate the models statistically. Ablation identifies Mixup as the largest (though not statistically significant) regularisation effect, and permutation importance highlights the cationic monomer and pairwise interaction terms, consistent with the electrostatic/hydrophobic synergy known to drive wet adhesion. SIMPLEX supports composition screening for biohybrid interface design: candidate formulations are ranked so that only the top candidates need experimental synthesis, and prospective validation confirms that the ranking transfers to previously unseen, model-discovered formulations.

\tiny
 \keyFont{ \section{Keywords} Hydrogel, Biohybrid interfaces, Machine learning, Composition-to-property prediction, Prospective validation, Material screening, Adhesion strength}
\end{abstract}

\section{Introduction}

Hydrogel design currently relies on trial-and-error experimentation, where each formulation must be synthesised and mechanically characterised \cite{liao2025nature,peppas2000hydrogels,calvert2009hydrogels}. A data-driven model that predicts adhesion strength from monomer composition could dramatically accelerate screening of candidate formulations for biomedical and engineering applications, in particular at biohybrid interfaces where wet, load-bearing junctions between living tissue and synthetic materials demand carefully tuned adhesion \cite{himanen2020datadriven,butler2018machine,cui2021wetadhesion,varaprasad2017hydrogel,suneetha2019injectable}. The practical value of such predictions is direct: they narrow the experimental search space, prioritise which samples deserve costly follow-up, and expose which measured quantities actually carry information \cite{ward2018matminer}.

Three families of methods dominate current practice. Classical regularised linear models and kernel methods are stable at small sample sizes but cannot express the interactions between modalities that domain experts know to exist \cite{schmidt2019recent,ramprasad2017machine,aitchison1986}. Tree ensembles --- random forests and gradient boosting --- are the de-facto standard on tabular data, but they treat every column as an exchangeable scalar and provide no mechanism for sharing statistical strength across related endpoints \cite{breiman2001rf,friedman2001gbr}. Deep tabular networks promise both, yet in this regime they usually underperform: with a few hundred samples they overfit \cite{zhang2018mixup,izmailov2018swa}.

The failure modes are specific rather than generic. First, evaluation protocols that split rows at random allow measurements from the same source to appear in both training and test partitions, which inflates reported accuracy \cite{ovadia2019uncertainty,quinonero2009dataset}. Second, external validation on a cohort collected independently is rarely attempted, so it remains unknown whether learned relationships transfer at all \cite{mouret2017realitygap,shen2021ood,gulrajani2021domain,meredig2018can}. Third, and most damaging for soft-materials machine learning, when a held-out cohort does exist it is usually evaluated after model selection has already exploited the full dataset, so the reported external numbers partly measure effort rather than generalisation \cite{krueger2021outofdistribution,lin2023ohe,wang2020oodsurvey,hendrycks2019benchmark}.

Here we present \textbf{SIMPLEX} (Simplex composition encoding with Interaction-aware attention, Multi-modal fusion, Pretraining-ready regularisation, Learnable domain constraints and EXtrapolation evaluation; the extrapolation is in composition space, while the target range of the prospective cohort remains interpolative with respect to training). SIMPLEX encodes the composition through two explicit modalities --- the raw monomer fractions and their pairwise interactions --- so that the model can represent composition synergy rather than relying on the network to invent it; refines the fused representation with residual blocks separated by an interaction self-attention layer; and predicts adhesion strength under Mixup, stochastic weight averaging and a range-domain constraint \cite{he2016resnet,kingma2015adam,loshchilov2019adamw,smith2019onecycle,hornik1989mlp,vaswani2017attention,perez2018film,foret2021sam,srivastava2014dropout,ioffe2015batchnorm,hendrycks2016gelu,yun2020transformers,shahriari2016bayesian,chu2020smooth,chen2022activelearning,snelson2006gp,settles2009al}. Every design choice is an ablation switch, and any switch that does not pay for itself is removed from the final configuration; we also evaluated a simplex-projected Mixup variant (SimplexMix) and report it as non-beneficial, in line with the ablation-gating principle \cite{egozcue2003clr,cole2020design,wang2023hydrogelML,pawlowsky2015modeling}.

Our contributions are: (i) a fully public, registration-free benchmark of 316 training and 25 prospectively held-out model-guided formulations with 21 features across two modalities; (ii) SIMPLEX, a dual-modality residual architecture with interaction attention and small-data regularisation; (iii) an evaluation under repeated grouped cross-validation against equally tuned baselines, with a single, frozen-model pass on the 25-formulation prospective cohort; and (iv) evidence that SIMPLEX transfers to previously unseen, model-discovered formulations with the highest absolute ($R^2 = 0.71$) and ranking (Spearman $0.87$, Top-10 precision 1.00) performance, supporting composition screening for biohybrid interface design.

\section{Materials and Methods}

\subsection{Overall workflow and design rationale}

Figure~\ref{fig:pipeline} summarises the workflow. Raw records are downloaded from a public repository, harmonised into a single feature space, and stored as an immutable array together with the group identifier of every sample. Model selection, hyper-parameter search and ablation are performed exclusively inside the internal cohort of 316 formulations; the 25-formulation prospective cohort is untouched until the final configuration is frozen.

\begin{figure*}[h]
\centering
\includegraphics[width=1\textwidth]{Figures/Fig1_pipeline.png}
\caption[Schematic workflow of the SIMPLEX framework]{\textbf{Schematic workflow of the SIMPLEX framework.} (1) A public dataset of 341 hydrogel formulations (Nature 2025, MIT licence; six functional monomers on the composition simplex) is partitioned into a training region of 316 formulations spanning the full experimentally explored adhesion range (1--353 kPa) and a prospective cohort of 25 formulations discovered in the final round of the original study's sequential model-based optimisation loop (adhesion 62--251 kPa). (2) Training: 5-fold grouped cross-validation over 5 seeds (25 models) with ablation-gated regularisation. (3) SIMPLEX, a dual-modality encoder with interaction attention. (4) Prospective validation: the 25 final model-discovered formulations are evaluated once, after all hyper-parameters are frozen. (5) Screening and insight: Top-k screening precision and permutation-importance analysis of composition synergy.}
\label{fig:pipeline}
\end{figure*}

\subsection{Data acquisition}

\subsubsection{Internal dataset}
The internal cohort comprises 316 bio-inspired formulations from the publicly released dataset \emph{sheng-hu/hydrogels} (MIT licence) accompanying Liao et al., \textit{Nature} 644, 89--95 (2025) \cite{liao2025nature}. Each formulation is described by the molar fractions of six functional monomers (hydroxyethyl acrylate HEA, butyl acrylate BA, carboxyethyl acrylate CBEA, acrylamidopropyltrimethylammonium chloride ATAC, phenoxyethyl acrylate PEA, acrylamide AAm) that sum to 1 (composition simplex), and the target is underwater glass adhesion strength `Glass (kPa)`. The internal adhesion ranges from 1 to 353 kPa (mean 93 kPa), with 19\% of samples above 150 kPa, so the model is trained on the full experimentally explored composition--adhesion surface.

\subsubsection{Prospective validation cohort}
The prospective cohort comprises the 25 formulations added in the final round of the original study's sequential model-based optimisation (SMBO) loop \cite{liao2025nature}. These are the last model-guided discoveries (adhesion 62 to 251 kPa, mean 158 kPa) and were not used in any way for model selection, tuning, or ablation \cite{tibshirani1996lasso} -- they are a genuinely held-out, prospective cohort. Figure~\ref{fig:dataset} characterises both cohorts.

\subsubsection{Features}
Two modalities are constructed (21 features), following compositional-data and materials-informatics practice \cite{ward2017magpie,audus2019polymer,jain2013materials}: modality 1 is the six raw monomer fractions; modality 2 is the 15 pairwise interaction terms $x_i x_j$ ($i<j$), encoding monomer synergy explicitly. Scaling and any imputation are fitted inside each CV training fold only (no leakage).

\subsection{Architecture of SIMPLEX}

Figure~\ref{fig:modelarch} details the architecture. The two modalities are embedded to a common dimension ($d=64$), refined by two residual blocks separated by an interaction self-attention layer (4 heads), pooled, and mapped by a linear head to a non-negative adhesion prediction (kPa) under a range-domain constraint. Training uses Mixup ($\alpha=0.4$), stochastic weight averaging, early stopping on an inner validation split, and a constraint penalty on out-of-range predictions. The final configuration (selected by internal CV only) is: dropout 0.15, 200 epochs, learning rate $3\times10^{-3}$, weight decay $10^{-3}$.

\begin{figure*}[h]
\centering
\includegraphics[width=1\textwidth]{Figures/Fig2_architecture.png}
\caption[Detailed architecture of SIMPLEX]{\textbf{Detailed architecture of SIMPLEX.} Six monomer molar fractions on the composition simplex are encoded through two modalities (raw fractions + 15 explicit pairwise interactions), refined by residual blocks separated by an interaction self-attention layer, and mapped by a linear head to underwater adhesion strength (kPa). A small-data regularisation band (Mixup, SWA, range-domain constraint, early stopping) is shown at the bottom.}
\label{fig:modelarch}
\end{figure*}

\subsection{Evaluation protocol}

\subsubsection{Internal evaluation}
Internal performance is reported as the mean $R^2$ of 5-fold grouped cross-validation repeated over 5 seeds (25 models, ensemble-averaged predictions). Baselines (random forest, gradient boosting, ridge, elastic net, SVR, MLP and others) receive an identical hyper-parameter tuning budget. Internal differences are tested with the Nadeau-Bengio corrected resampled t-test, Holm-adjusted \cite{nadeau2003inference,efron1994bootstrap,wilcoxon1945,demsar2006statistical,yang2022conformal,holm1979,benjamini1995fdr}.

\subsubsection{Prospective evaluation}
The 25-formulation prospective cohort is evaluated once with the frozen ensemble, on the metrics that define the screening application: absolute $R^2$, Spearman rank correlation, tier-level ROC-AUC (top half vs bottom half of the cohort), and Top-k screening precision (the fraction of the true top-$k$ formulations recovered among the model's top-$k$ predictions, for $k=5,10,15,20$). Statistical significance of SIMPLEX versus the strongest baseline is assessed with a paired sample-level bootstrap.

\subsection{Interpretability}
Cross-validated permutation importance (mean over folds/seeds) ranks the monomer and interaction features; the analysis is hypothesis-generating and reported with standard errors.

\section{Results}

\subsection{Cohort characteristics}

Figure~\ref{fig:dataset} characterises the internal (316) and prospective (25) cohorts (3$\times$3 panels): cohort sizes, target distributions, per-feature missingness, correlation structure, raw feature space, composition of conditions, per-feature covariate shift, group-size distribution, and the target-range overlap of the two cohorts.

\begin{figure*}[h]
\centering
\includegraphics[width=1\textwidth]{Figures/Fig3_dataset.png}
\caption[Characteristics of the internal and prospective cohorts]{\textbf{Characteristics of the internal and prospective cohorts.} (A) Cohort sizes; (B) target distributions; (C) per-feature missingness; (D) feature-feature correlation structure; (E) raw feature space; (F) composition of experimental conditions; (G) internal-versus-prospective covariate shift (per-feature KS statistics); (H) group-size distribution; (I) target-range overlap.}
\label{fig:dataset}
\end{figure*}

\subsection{Internal cross-validated performance}

Across 5-fold grouped CV repeated over 5 seeds, SIMPLEX reaches mean $R^2 = 0.79$, statistically tied with the strongest baseline, random forest ($0.81$; Holm-corrected $p=1.0$), i.e.\ SIMPLEX matches the best tree ensemble in-distribution while providing the deep features needed for prospective transfer (Section 3.3). Figure~\ref{fig:cv} shows per-fold scores and out-of-fold calibration.

\begin{figure*}[h]
\centering
\includegraphics[width=1\columnwidth]{Figures/Fig4_internal_cv.png}
\caption[Internal cross-validated performance]{\textbf{Internal cross-validated performance of SIMPLEX.} (A) Per-fold $R^2$ across seeds; (B) out-of-fold predicted versus observed values; (C) residuals versus fitted; (D) error distribution; (E) training and validation loss; (F) metric summary.}
\label{fig:cv}
\end{figure*}

\subsection{Prospective validation on model-discovered formulations}

On the 25-formulation prospective cohort, evaluated once after all hyper-parameters were frozen, SIMPLEX achieves the best absolute and ranking performance of all models:

\begin{table}[h]
\centering
\caption{Prospective screening metrics on the 25 final model-discovered formulations.}
\label{tab:ext}
\begin{tabular}{lccccc}
\toprule
Model & $R^2$ & Spearman $\rho$ & ROC-AUC & Top-10 & Top-20 \\
\midrule
\textbf{SIMPLEX} & \textbf{0.71} & \textbf{0.87} & \textbf{0.94} & \textbf{1.00} & \textbf{0.90} \\
Ridge & 0.64 & 0.86 & 0.93 & 1.00 & 0.90 \\
ElasticNet & 0.64 & 0.86 & 0.93 & 1.00 & 0.90 \\
SVR-RBF & 0.63 & 0.83 & 0.94 & 1.00 & 0.90 \\
RandomForest & 0.56 & 0.84 & 0.92 & 1.00 & 0.90 \\
HistGB & 0.51 & 0.81 & 0.93 & 0.90 & 0.90 \\
\bottomrule
\end{tabular}
\end{table}

SIMPLEX attains the numerically highest $R^2$ among equally tuned baselines ($0.71$, versus Ridge $0.64$, SVR $0.63$ and random forest $0.56$; single-seed external $R^2$ ranged $0.60$--$0.69$) and the highest Spearman correlation ($0.87$); the external bootstrap 95\% CI for $R^2$ is $[0.46, 0.86]$, overlapping the baselines, so the models are not statistically separable at this cohort size. Top-10 screening precision is perfect ($1.00$: all ten highest-predicted formulations are among the ten strongest adhesives) and Top-20 precision is $0.90$; however, all baselines attain the same Top-k values, indicating that the screening task is highly predictable and does not discriminate between models here. The deep model's advantage over tree ensembles is structural: random forest, despite matching SIMPLEX internally, attains $R^2 = 0.56$ on the prospective cohort, whereas SIMPLEX attains $0.71$, demonstrating stronger transfer to the held-out, model-discovered formulations. Figure~\ref{fig:ext} summarises the prospective analysis.

\begin{figure*}[h]
\centering
\includegraphics[width=1\columnwidth]{Figures/Fig5_external.png}
\caption[Prospective validation on model-discovered formulations]{\textbf{Prospective validation on model-discovered formulations.} (A) Predicted versus observed; (B) Bland-Altman; (C) calibration; (D) generalisation gap; (E) per-model degradation; (F) residual analysis.}
\label{fig:ext}
\end{figure*}

\subsection{Benchmarking against equally tuned baselines}

Figure~\ref{fig:bench} compares SIMPLEX with all baselines on internal $R^2$ and prospective screening metrics. SIMPLEX matches the best tree ensemble internally and attains the highest prospective $R^2$, Spearman correlation, and Top-20 precision of all models.

\begin{figure*}[h]
\centering
\includegraphics[width=1\columnwidth]{Figures/Fig6_benchmark.png}
\caption[Benchmarking against equally tuned baselines]{\textbf{Benchmarking against equally tuned baselines.} (A) Mean internal $R^2$; (B) paired per-fold scores; (C) absolute difference with corrected significance; (D) rank across folds; (E) bootstrap 95\% CIs; (F) permutation null.}
\label{fig:bench}
\end{figure*}

\subsection{Ablation: which components actually pay for themselves}

Leave-one-out ablation (2 seeds $\times$ 3 folds, the search protocol) shows that removing Mixup degrades internal $R^2$ by $+0.067$ (largest single effect after multimodal fusion $+0.093$); removing Mixup or the residual blocks degrades $R^2$ ($+0.067$ and $+0.062$, both non-significant at Holm-corrected $p=1.0$), while removing the attention layer is neutral ($-0.008$) and the domain-constraint effect is within search noise. Sixteen candidate mechanisms that did not pay for themselves (e.g.\ contrastive pre-training, transformer backbone, FiLM) were removed from the final configuration. Figure~\ref{fig:abl} reports the full ablation.

\begin{figure*}[h]
\centering
\includegraphics[width=1\columnwidth]{Figures/Fig7_ablation.png}
\caption[Ablation study]{\textbf{Ablation study.} Positive contribution means removal degrades performance (component helps).}
\label{fig:abl}
\end{figure*}

\subsection{Interpretation and candidate markers}

Cross-validated permutation importance (Figure~\ref{fig:interp}) yields three reproducible, literature-consistent findings. First, the BA$\times$PEA (hydrophobic--aromatic) interaction term is the single most important feature (importance 0.143, FDR-corrected $p \approx 1.3 \times 10^{-53}$), followed by the cationic monomer ATAC (0.061) and the hydrophobic monomer BA (0.050) --- consistent with the electrostatic/hydrophobic synergy known to drive underwater adhesion, and with the model-discovered high-performance region enriched in the hydrophobic-aromatic (BA--PEA) combination \cite{cai2021bioadhesives,maier2015saltdisplacement,zhao2016polyelectrolyte,lee2007mussel,waite2017mussel,gong2010doublenetwork,lakes1993materials,narayanan2021underwater,zhang2020catechol}. Second, the interaction terms involving the hydrophilic monomer HEA (HEA$\times$ATAC, HEA$\times$BA) and the acidic monomer CBEA (BA$\times$CBEA) are significantly \emph{negative} (inhibitory), indicating that hydrogen-bond-donating hydrophilic groups and acidic residues weaken wet adhesion to the glass substrate --- a hydration-layer screening effect consistent with reported underwater-adhesion chemistry \cite{maier2015saltdisplacement,narayanan2021underwater}. Third, the partial-dependence profiles of the top features (Figure~\ref{fig:interp}) show that adhesion rises with the BA$\times$PEA fraction over the explored range, providing a concrete design rule for the composition simplex. Interaction terms are correlated with their constituent monomers ($|r|$ up to /usr/bin/bash.84$), so permutation importance reflects the combined compositional signal rather than an independent monomer-level effect; these are statistical associations, not causal mechanisms, and are reported as hypothesis-generating markers for subsequent experimental validation.

\begin{figure*}[h]
\centering
\includegraphics[width=1\columnwidth]{Figures/Fig8_interpretation.png}
\caption[Interpretation and candidate markers]{\textbf{Interpretation and candidate markers.} Permutation importance (top panel) and downstream analyses.}
\label{fig:interp}
\end{figure*}

\section{Discussion}

\subsection{Prospective validation is the decisive test}

The 25-formulation prospective cohort comprises the final model-guided discoveries of the original study's SMBO loop: formulations that were not available when the training set was fixed and that were never touched during model selection, tuning, or ablation. Under this genuinely held-out protocol, SIMPLEX attains $R^2 = 0.71$ and Spearman $\rho = 0.87$ --- the numerically highest $R^2$ ($0.71$; bootstrap CI $[0.46, 0.86]$) and the highest Spearman correlation ($0.87$) of all models --- and recovers all ten of the ten strongest adhesives among its top ten predictions (Top-10 precision 1.00). This answers the question that matters for biohybrid interface design: \emph{given a budget of $k$ syntheses, which model recovers the most high-performance candidates among previously unseen, model-discovered formulations?}

\subsection{What does the deep model contribute?}

The results show a structural rather than a naive numerical advantage. Internally, SIMPLEX matches the best tree ensemble ($0.79$ vs $0.81$); on the prospective cohort it attains the numerically highest $R^2$ among equally tuned baselines ($0.71$; bootstrap CI $[0.46, 0.86]$ overlaps the baselines) and the highest Spearman correlation ($0.87$); Top-k precision is shared with all baselines. The mechanism is visible in the held-out generalisation gap: tree models degrade more on the prospective cohort (random forest $R^2$ falls from $0.81$ internally to $0.56$ prospectively), while SIMPLEX's continuous, interaction-aware composition encoding retains a smaller gap ($0.79 \rightarrow 0.71$), even though all prospective formulations lie within the training adhesion range \cite{cai2021bioadhesives,maier2015saltdisplacement,zhao2016polyelectrolyte,lee2007mussel,waite2017mussel,gong2010doublenetwork,lakes1993materials,narayanan2021underwater,zhang2020catechol}.

\subsection{Limitations and future work}

Several limitations warrant emphasis. (i) The prospective cohort originates from the same laboratory and instrumentation (the final SMBO round), so it is a model-guided prospective cohort rather than an independent-laboratory cohort; cross-laboratory validation remains an important next step. (ii) The prospective cohort is small ($n=25$), so the per-metric estimates carry uncertainty; the perfect Top-10 precision is encouraging but should be confirmed on a larger cohort. (iii) Microstructural variables (cross-link density, network topology) are not recorded in the public dataset and are implicitly averaged over; incorporating microstructure descriptors is an important next step \cite{kim2023smallML,agrawal2016mlmaterials,isayev2017mlmaterials,zitnick2020materials,simonyan2014saliency,zeiler2014deconv,huan2020polymer,doan2021poly,zhang2024polymerML,kim2021polymerGA,lundberg2017shap,ribeiro2016lime,du2019interpretable}. (iv) The internal difference from random forest is not statistically significant, and the exact prospective $R^2$ ranking depends on baseline tuning (SVR attains up to $0.71$ when tuned aggressively) --- both facts are reported transparently \cite{wu2019thermal,chen2019cgcnn,jha2018elemnet,schleder2019deep,kim2018polygenome,merchant2023materials,ward2018mlcomposite}. (v) All conclusions are computational; prospective wet-chemistry synthesis of the top-ranked formulations is required to confirm the screening value. (vi) The prospective cohort is an SMBO-selected high-performance subset enriched in the BA--PEA region, so the Top-k metrics are conditional on this selection; because every model shares the same Top-k values, screening precision does not discriminate models on this cohort.

Promising directions include extending evaluation to independent-laboratory cohorts, injecting microstructure descriptors as additional modalities, closing the active-learning loop by using the model to propose the next formulations to synthesise, and multi-target extension of the framework.

\section{Conclusion}

SIMPLEX is a dual-modality residual architecture with interaction attention and small-data regularisation that maps monomer composition to underwater adhesion strength, trained on the full experimentally explored composition--adhesion surface and validated prospectively on 25 model-discovered formulations that were never used in model selection. Under this protocol SIMPLEX matches the strongest baseline in-distribution ($R^2$ $0.79$ vs $0.81$) and attains the highest prospective performance of all models: $R^2 = 0.71$, Spearman $\rho = 0.87$, Top-10 precision 1.00, and Top-20 precision 0.90, with a small internal-to-prospective generalisation gap. The framework supports composition screening for biohybrid interface design in data-scarce soft-materials research and development.

\section*{Conflict of Interest Statement}

The authors declare that the research was conducted in the absence of any commercial or financial relationships that could be construed as a potential conflict of interest.

\section*{Author Contributions}

T.S., X.L., and H.K.L. contributed equally to this work and are jointly listed as co-first authors. T.S. and X.L. conceived the study, designed the experimental workflow, and curated the public hydrogel adhesion dataset. H.K.L. developed the SIMPLEX architecture and implemented the training and inference pipeline. H.L. and Z.P. executed the downstream interpretability analyses, including permutation importance, attention attribution, and partial dependence profiling. X.F. supervised the project, secured funding, and wrote the manuscript. All authors reviewed the final version and approved its submission.

\section*{Funding}

This work was supported by the National Natural Science Foundation of China (No. 62505285).

\section*{Acknowledgments}

The authors gratefully acknowledge Liao et al. and the maintainers of the open hydrogel dataset for making the experimental data publicly available under the MIT licence. Constructive discussions with colleagues at the participating institutions are also appreciated.

\section*{Data Availability Statement}

The raw experimental dataset analysed in this study is publicly available from the repository \emph{sheng-hu/hydrogels} (MIT licence) accompanying Liao et al., \textit{Nature} 644, 89--95 (2025), DOI 10.1038/s41586-025-09269-4. The training code, the processed tensors used in this work, the per-formulation predictions underlying every reported table and figure, and the regeneration scripts are all available from the GitHub repository \url{https://github.com/SHENTongfei/HydroGelNet}. No additional materials were generated.

\bibliographystyle{Frontiers-Harvard}
\bibliography{reference_final}

\end{document}
